\section{Conclusion}
\label{sec:conclusion}

This paper studied the allocation of expensive true evaluations in the grey wolf
optimizer. The starting observation is structural: GWO's search direction is
produced by three positions, yet the canonical algorithm evaluates its whole
population every generation. We proposed RASE-GWO, which separates a wolf's
virtual search position from its verified position, keeps a candidate pool that
is an order of magnitude larger than the per-round evaluation quota, and assigns
that quota by a surrogate-free geometric rule that combines consensus refinement
with minority-leader probing. Only verified positions may enter the
$\alpha/\beta/\delta$ hierarchy, which is what makes sparse evaluation sound for
a leader-driven algorithm. A small Q-learner adapts the split between the two
groups without knowledge of the problem landscape.

The empirical study on CEC2017 at $D=50$ holds the true-evaluation budget, the
number of generations and the initial populations fixed, so every difference
between the compared algorithms is attributable to candidate selection. Under
that protocol the proposed mechanism improves on canonical GWO, and the ablation
isolates the contribution of the geometric screening, the over-sampled candidate
pool and the verification decoupling. Equally important, the study delimits the
framework: reducing the per-round quota without over-sampling does not pay off
for GWO, crowd-level diversity scoring is ineffective in a leader-driven
algorithm, and shrinking the population is not a substitute for screening a
larger candidate pool.

Future work follows three directions: extending the screening to multiple
objectives and to constrained problems, replacing the tabular Q-learner with a
function approximation that can condition on measurable landscape features, and
applying the framework to expensive engineering simulators, where position
updates are negligible compared with a single objective evaluation.
